% ======================================================================
% DISCUSSION — concise, quantitative, no marketing superlatives.
% ======================================================================
\section*{Discussion}

We have shown that an ordered magnetic cilia layer read out by distributed
tri-axis sensors provides a route to embodied perception of subtle
mechanical stimuli. Relative to conventional planar magnetic skins, the
cilia geometry enhances the weak-stimulus response by
\dataneeded{$G_{\mathrm{cilia}}$}-fold while retaining the contactless,
drift-tolerant, multi-axis character of magnetic transduction. The local
$2\times2$ unit generalizes to unseen contact positions, and the same
decoding scales to a 44-node hand with explicit artifact
suppression---addressing the central risk that whole-hand maps reflect
posture or field artifacts rather than true contact.

Several limitations bound these claims. First, ground-truth forces are
established with an ATI~Nano17, whose resolution does not certify sub-mN or
$\mu$N-level truth; the bubble-touch result is therefore an
ultralight-contact demonstration, not a detection limit. Second, the
airflow maps report airflow-induced tactile responses, not a resolved flow
field, which would require PIV or multi-point anemometry. Third, pulse
waveforms are validated against ECG in timing only; morphological features
such as the dicrotic notch require a pressure or PPG reference for
confirmation. Finally, in the absence of feedback experiments we claim
grasp-state discrimination and distributed contact readout, not closed-loop
control.

The comparison with vision-based tactile sensing clarifies where the
platform sits. Vision-based sensors reach sub-millimeter local resolution
but require an internal camera, illumination and per-contact computation,
which constrain thickness, curvature and integration density. The magnetic
cilia tactile skin makes the opposite trade: it accepts coarser local
resolution in exchange for a camera-free, low-profile, three-axis form
factor that covers a full hand and remains sensitive to weak dynamic
disturbances such as airflow and pulse micro-vibrations. Future work
includes per-node force calibration across the whole hand, larger-cohort
pulse studies with a PPG reference, and integration of the tactile maps
into manipulation feedback loops.
